\SourceLeaf{372}%
Es ist kaum begreiflich, wie er sich auf eine solche Weise selbst täuschen konnte. Er macht sich sogar selbst den Einwurf, dass die Summe der Stücke \(Sp,pp'\) u.\,s.\,w., wenn die Stücke immer kleiner und kleiner werden, doch, ungeachtet ihre Anzahl ohne Aufhören zunehme, nicht über eine gewisse Grenze hinauswachsen könnte, und meint diesen Einwurf damit zu heben, dass jene Stücke, auch wenn sie immer kleiner und kleiner werden, doch immer grösser bleiben, als \textit{eine angebliche Grösse}; nemlich jene Stücke sind Katheten von rechtwinkligen Dreiecken, und folglich immer grösser als der Unterschied zwischen Hypotenuse und der andern Kathete. Fast scheint es, dass eine grammatische Zweideutigkeit den Verf. irre geleitet hat, nemlich der zwiefache Sinn des Artikels \textit{eine} angebliche Grösse. Der Schluss des Verf. würde nur dann richtig sein, wenn sich zeigen liesse, dass die Stücke \(Sp,pp'\) u.\,s.\,w. immer grösser bleiben. als \textit{eine bestimmte} angebliche Grösse, z.\,B. als der Unterschied zwischen der Hypotenuse \(pT\) und der Kathete \(ST\). Aber das lässt sich nicht beweisen, sondern nur, dass jedes Stück immer grösser bleibt, als eine angebliche Grösse, die aber selbst für jedes Stück eine andere ist, nemlich \(Sp\) grösser als der Unterschied zwischen \(pT\) und \(ST\), ferner \(pp'\) grösser als der Unterschied zwischen \(qp'\) und \(qp\) u.\,s.\,w. Hiemit verschwindet nun aber die ganze Kraft des Beweises.

Auf dieselbe Art, wie er seinen Beweis führen zu können geglaubt hat, könnte er auch beweisen, dass in einem ebnen Dreiecke \(ABC\), worin \(B\) ein rechter Winkel ist, \(C\) nicht spitz sein könne; er brauchte nur aus \(B\) ein Perpendikel \(BD\) auf die Hypotenuse \(AC\) zu fällen, dann wieder das Perpendikel \(DE\) auf \(AB\) und so ohne Aufhören die Perpendikel \(EF,FG,GH\) u.\,s.\,w. wechselsweise auf \(AC\) und \(AB\). Die Stücke \(CD,DF,FH\) u.\,s.\,w. sind immer grösser als der angebliche Unterschied zwischenHypotenuse und einer Kathete \textit{desjenigen} rechtwinkligen Dreiecks, worin jede der Reihe nach die andere Kathete ist, demungeachtet erschöpft ihre Summe offenbar die Hypotenuse \(AC\) nie, so gross auch ihre Anzahl genommen wird.

Wir müssten fast bedauern, bei so bekannten und leichten Dingen so lange verweilt zu haben, wenn nicht diese Schrift, deren Verf. es übrigens wirklich um Wahrheit zu thun zu sein scheint, durch die Art wie sie schon vor ihrer Erscheinung in öffentlichen Blättern angekündigt wurde, eine mehr als gewöhnliche Aufmerksamkeit auf sich gezogen hätte. Wir bemerken daher hier nur noch, dass der Verf. nachher auf eine ganz ähnliche, und daher eben so nichtige Art bewei%
